\documentclass[11pt,a4paper]{article}
% preamble.tex — estilo común de todos los apuntes Froth.
% Cada apunte hace % preamble.tex — estilo común de todos los apuntes Froth.
% Cada apunte hace % preamble.tex — estilo común de todos los apuntes Froth.
% Cada apunte hace \input{preamble} justo después de \documentclass.
\usepackage[margin=2.4cm]{geometry}
\usepackage{xcolor}
\definecolor{litblue}{RGB}{31,79,121}
\definecolor{litgreen}{RGB}{27,94,32}
\definecolor{litorange}{RGB}{230,81,0}
\usepackage{parskip}
\usepackage{titlesec}
\titleformat{\section}{\Large\bfseries\color{litblue}}{\thesection}{0.6em}{}
\titleformat{\subsection}{\large\bfseries\color{litblue}}{\thesubsection}{0.6em}{}
\usepackage{tcolorbox}
\usepackage{fancyhdr}
\pagestyle{fancy}
\fancyhf{}
\fancyhead[L]{\small\color{litblue}Froth — Apuntes de ML}
\fancyhead[R]{\small\thepage}
\renewcommand{\headrulewidth}{0.4pt}
\usepackage[colorlinks=true,linkcolor=litblue,urlcolor=litblue]{hyperref}

% Cabecera de cada apunte: \cabecera{numero}{titulo}
\newcommand{\cabecera}[2]{%
  \begin{tcolorbox}[colback=litblue,coltext=white,colframe=litblue,arc=2mm]
    {\Large\bfseries Apunte #1 — #2}\\[3pt]
    {\small Proyecto Froth · aprender Machine Learning construyendo}
  \end{tcolorbox}\vspace{0.8em}}

% Cajas temáticas
\newtcolorbox{concepto}[1]{colback=blue!4,colframe=litblue,title={Concepto: #1},arc=1.5mm}
\newtcolorbox{practica}{colback=green!5,colframe=litgreen,title={Pruébalo tú (teclado en mano)},arc=1.5mm}
\newtcolorbox{ojo}{colback=orange!8,colframe=litorange,title={Ojo / error típico},arc=1.5mm}
\newtcolorbox{resumen}{colback=gray!8,colframe=gray!60,title={En una frase},arc=1.5mm}

\newcommand{\codigo}[1]{\texttt{#1}}
 justo después de \documentclass.
\usepackage[margin=2.4cm]{geometry}
\usepackage{xcolor}
\definecolor{litblue}{RGB}{31,79,121}
\definecolor{litgreen}{RGB}{27,94,32}
\definecolor{litorange}{RGB}{230,81,0}
\usepackage{parskip}
\usepackage{titlesec}
\titleformat{\section}{\Large\bfseries\color{litblue}}{\thesection}{0.6em}{}
\titleformat{\subsection}{\large\bfseries\color{litblue}}{\thesubsection}{0.6em}{}
\usepackage{tcolorbox}
\usepackage{fancyhdr}
\pagestyle{fancy}
\fancyhf{}
\fancyhead[L]{\small\color{litblue}Froth — Apuntes de ML}
\fancyhead[R]{\small\thepage}
\renewcommand{\headrulewidth}{0.4pt}
\usepackage[colorlinks=true,linkcolor=litblue,urlcolor=litblue]{hyperref}

% Cabecera de cada apunte: \cabecera{numero}{titulo}
\newcommand{\cabecera}[2]{%
  \begin{tcolorbox}[colback=litblue,coltext=white,colframe=litblue,arc=2mm]
    {\Large\bfseries Apunte #1 — #2}\\[3pt]
    {\small Proyecto Froth · aprender Machine Learning construyendo}
  \end{tcolorbox}\vspace{0.8em}}

% Cajas temáticas
\newtcolorbox{concepto}[1]{colback=blue!4,colframe=litblue,title={Concepto: #1},arc=1.5mm}
\newtcolorbox{practica}{colback=green!5,colframe=litgreen,title={Pruébalo tú (teclado en mano)},arc=1.5mm}
\newtcolorbox{ojo}{colback=orange!8,colframe=litorange,title={Ojo / error típico},arc=1.5mm}
\newtcolorbox{resumen}{colback=gray!8,colframe=gray!60,title={En una frase},arc=1.5mm}

\newcommand{\codigo}[1]{\texttt{#1}}
 justo después de \documentclass.
\usepackage[margin=2.4cm]{geometry}
\usepackage{xcolor}
\definecolor{litblue}{RGB}{31,79,121}
\definecolor{litgreen}{RGB}{27,94,32}
\definecolor{litorange}{RGB}{230,81,0}
\usepackage{parskip}
\usepackage{titlesec}
\titleformat{\section}{\Large\bfseries\color{litblue}}{\thesection}{0.6em}{}
\titleformat{\subsection}{\large\bfseries\color{litblue}}{\thesubsection}{0.6em}{}
\usepackage{tcolorbox}
\usepackage{fancyhdr}
\pagestyle{fancy}
\fancyhf{}
\fancyhead[L]{\small\color{litblue}Froth — Apuntes de ML}
\fancyhead[R]{\small\thepage}
\renewcommand{\headrulewidth}{0.4pt}
\usepackage[colorlinks=true,linkcolor=litblue,urlcolor=litblue]{hyperref}

% Cabecera de cada apunte: \cabecera{numero}{titulo}
\newcommand{\cabecera}[2]{%
  \begin{tcolorbox}[colback=litblue,coltext=white,colframe=litblue,arc=2mm]
    {\Large\bfseries Apunte #1 — #2}\\[3pt]
    {\small Proyecto Froth · aprender Machine Learning construyendo}
  \end{tcolorbox}\vspace{0.8em}}

% Cajas temáticas
\newtcolorbox{concepto}[1]{colback=blue!4,colframe=litblue,title={Concepto: #1},arc=1.5mm}
\newtcolorbox{practica}{colback=green!5,colframe=litgreen,title={Pruébalo tú (teclado en mano)},arc=1.5mm}
\newtcolorbox{ojo}{colback=orange!8,colframe=litorange,title={Ojo / error típico},arc=1.5mm}
\newtcolorbox{resumen}{colback=gray!8,colframe=gray!60,title={En una frase},arc=1.5mm}

\newcommand{\codigo}[1]{\texttt{#1}}

\begin{document}
\cabecera{03}{Cuando internet dice que no: SSL, códigos de error y API keys}

Hoy chocamos con dos errores distintos seguidos. Vale oro saber distinguirlos, porque uno
era \emph{de tu máquina} y el otro \emph{del servidor} — y se arreglan de formas opuestas.

\section{Error 1: el candado (SSL / certificados)}

\begin{concepto}{certificados y HTTPS}
Cuando tu PC se conecta a \codigo{https://api.openalex.org}, antes de mandar nada verifica
el \textbf{certificado} del servidor: su carnet de identidad digital, firmado por
autoridades de confianza. Así sabes que hablas con OpenAlex de verdad y no con un impostor.

Python trae su propia lista de autoridades de confianza. En tu red (corporativa/con
antivirus que inspecciona el tráfico), el certificado que llega no estaba en esa lista
$\rightarrow$ Python se negó: \codigo{CERTIFICATE\_VERIFY\_FAILED}. Tu navegador sí
funcionaba porque usa la lista de \emph{Windows}, que sí conoce ese certificado.

\textbf{Arreglo:} la librería \codigo{truststore} le dice a Python ``usa la lista de
Windows''. La activamos una sola vez en \codigo{froth/\_\_init\_\_.py} y vale para todo
el proyecto.
\end{concepto}

\section{Error 2: la puerta cerrada (código 503)}

\begin{concepto}{códigos de estado HTTP}
Toda respuesta HTTP trae un número de 3 cifras que resume cómo fue:
\begin{itemize}
  \item \textbf{2xx} --- bien (\codigo{200 OK}).
  \item \textbf{4xx} --- el error es \emph{tuyo}: \codigo{404} no existe, \codigo{401} sin permiso,
        \codigo{429} pides demasiado rápido.
  \item \textbf{5xx} --- el error es \emph{del servidor}: \codigo{500} se rompió por dentro,
        \codigo{503} ``no disponible'' (saturado o en mantenimiento).
\end{itemize}
Nuestro \codigo{503} venía con un mensaje claro: ``búsqueda anónima limitada por carga;
usa una API key gratuita''. No había nada que arreglar en el código: era su puerta, cerrada
para anónimos.
\end{concepto}

\begin{ojo}
Regla mental para siempre: \textbf{SSL/certificados $\rightarrow$ problema en tu lado}
(red, Python, certificados). \textbf{5xx $\rightarrow$ problema en el lado de ellos}
(esperar, identificarse o cambiar de servicio). Reintentar en bucle sin entender cuál de
los dos es, es perder el tiempo.
\end{ojo}

\section{La API key: tu carnet VIP (y por qué es un secreto)}

\begin{concepto}{API key}
Una \textbf{API key} es una cadena de texto que identifica tus peticiones ante el
servicio. Con ella OpenAlex te da acceso estable (1 USD/día de uso gratuito, unas 1000
búsquedas — de sobra). Se envía como un parámetro más: \codigo{api\_key=...}.

Como te identifica \emph{a ti}, es un \textbf{secreto}: quien la tenga gasta tu cuota en
tu nombre. Por eso:
\begin{itemize}
  \item \textbf{Nunca} se escribe dentro del código fuente.
  \item Vive en el archivo \codigo{.openalex\_key} en la raíz del proyecto (una sola línea).
  \item Ese archivo está en \codigo{.gitignore}: la lista de cosas que Git (control de
        versiones) tiene prohibido subir a ningún sitio. Si un día publicas el código en
        GitHub, tu clave no viaja con él.
\end{itemize}
\codigo{config.py} la lee al arrancar (primero busca la variable de entorno
\codigo{OPENALEX\_API\_KEY}; si no, el archivo) y \codigo{pull.py} la añade a cada petición.
\end{concepto}

\begin{practica}
Comprueba tú mismo que la clave está donde debe y fuera del alcance de Git:
\begin{enumerate}
  \item Abre \codigo{.gitignore} (con el Bloc de notas) y localiza la línea \codigo{.openalex\_key}.
  \item En PowerShell, dentro de la carpeta del proyecto:
        \codigo{Get-Content .openalex\_key} \quad (verás tu clave — solo tú).
\end{enumerate}
\end{practica}

\begin{resumen}
SSL falló en tu máquina (arreglado con truststore); el 503 era el servidor saturado
rechazando anónimos (arreglado con API key); y una API key es un secreto que vive fuera
del código y fuera de Git.
\end{resumen}

\end{document}
